\begin{theorem}[A geometrically simply connected surface test]
\label{thm:B-simply-connected-surface-gate}
\status{Proved strict intrinsic test-class equivalence; restricted Tate unresolved}
Fix a finite field \(k\) of characteristic \(p\), a prime
\(\ell\ne p\), an integer \(N\geq1\), and a real number
\(\varepsilon>0\).  Divisor Tate for all smooth projective
geometrically integral surfaces over \(k\) is equivalent to
divisor Tate for just those surfaces \(Y/k\) satisfying
\[
 \pi_1^{\et}(Y_{\overline k})=1,
 \qquad H^1(Y,\mathcal O_Y)=0,
 \qquad \omega_Y\text{ very ample},
\]
and, simultaneously,
\[
 K_Y^2\geq N,\qquad p_g(Y)\geq N,\qquad c_2(Y)>0,
 \qquad \left|\frac{K_Y^2}{c_2(Y)}-1\right|<\varepsilon.
\]
No extension of the fixed field \(k\) is needed in the reduction.
Consequently the universal prime-to-characteristic torsor and
relative-Brauer finiteness assertions of
Theorem~\ref{thm:B-Brauer-prime-uniform} and
Corollary~\ref{cor:B-torsors-prime-uniform} have this smaller
intrinsic surface test class, when quantified over finite constant
extensions as before.  The numerical conditions are inherited from
Corollary~\ref{cor:B-canonical-surface-gate}, not independent
additional reductions.
\end{theorem}
\begin{proof}
The implication from all surfaces to the indicated subclass is
immediate.  For the converse start with a smooth projective
geometrically integral surface \(X/k\).  The already audited
proof of Yuan's Theorem~1.6 reduces its divisor Tate assertion,
over this same \(k\), to BSD for a semistable abelian variety
\(A/K\), where \(K=k(t)\), whose N\'eron Hodge bundle over
\(\mathbf P^1_k\) is ample.  If \(A=0\), the original Tate
assertion is already known.  Otherwise Yuan's Theorem~4.5,
applied to \(A^\vee\), gives a projective flat regular model
\(\psi:P\to\mathbf P^1_k\), with generic fibre \(A^\vee\),
and
\[
 R^1\psi_*\mathcal O_P\simeq
       \operatorname{Lie}(\mathcal A/\mathbf P^1_k),
 \qquad H^1(P,\mathcal O_P)=0.
\]
Here \(\mathcal A\) is the N\'eron model of \(A\); the
vanishing follows from ampleness of the dual bundle and Leray
over \(\mathbf P^1\)
\citep[Theorem~4.5 and Section~4.2, author pp.~33--36]{DYuan2021Hodge}.
The generic zero section extends over every point of the regular
base curve by properness, giving a global section and a
\(k\)-rational point.  Since \(k\) is perfect, \(P\) is smooth
over \(k\).

\emph{The geometric fundamental group of \(P\).}
Shrink to a nonempty open \(U\subset\mathbf P^1_k\) on which
the model is the abelian scheme extending its generic fibre.
The proper smooth homotopy sequence gives a right-exact sequence
\[
 \pi_1(P_{\overline\eta})\longrightarrow
 \pi_1(P_{U,\overline k})\longrightarrow
 \pi_1(U_{\overline k})\longrightarrow1,
\]
split on the right by the section
\citep[Expos\'e~X, Theorem~1.3 and Corollary~1.4,
pp.~202--203]{c_sga1_2003}.
No left injectivity is asserted.  Normality makes
\(\pi_1(P_{U,\overline k})\to\pi_1(P_{\overline k})\)
surjective.  The section subgroup dies in the latter group,
because it factors through \(\pi_1(\mathbf P^1_{\overline k})=1\).
Thus
\(G=\pi_1(P_{\overline k})\) is a quotient of the abelian
group \(\pi_1(P_{\overline\eta})\), and is itself abelian
\citep[Lemma~58.10.7]{BStacksFundamentalGroups};
\citep[Expos\'e~XI, Section~2 and Theorem~2.1,
pp.~221--222]{c_sga1_2003}.

The vanishing of \(H^1(P_{\overline k},\mathcal O)\) makes
the tangent space of the connected Picard scheme zero, so that
connected Picard scheme is trivial.  For every \(m\) prime to
\(p\), Kummer therefore gives
\[
 H^1(P_{\overline k},\mu_m)
       \simeq\operatorname{Pic}(P_{\overline k})[m]
       =\operatorname{NS}(P_{\overline k})[m].
\]
The torsion subgroup of N\'eron--Severi is one finite group;
hence the entire prime-to-\(p\) character group of \(G\) is
finite.  Artin--Schreier gives
\(H^1(P_{\overline k},\Z/p)=0\), since \(F-1\) is
surjective on \(\overline k\) and \(H^1(\mathcal O)=0\).
An abelian profinite group with a nonzero pro-\(p\) part has a
quotient \(\Z/p\).  That part is therefore zero.  Characters
separate points of an abelian profinite group, so \(G\) is
finite, of order prime to \(p\)
\citep[Expos\'e~XI, Proposition~6.4, Corollary~6.6 and its
following paragraph, pp.~234--235; Propositions~6.7--6.8 and
Corollary~6.10, pp.~236--237]{c_sga1_2003}.

\emph{A cover over the original field.}
Use the rational point on the zero section to split the arithmetic
short exact sequence
\[
 1\longrightarrow G\longrightarrow\pi_1(P)
      \longrightarrow\Gal(\overline k/k)\longrightarrow1.
\]
This arithmetic sequence, unlike the general homotopy sequence
above, is left exact as well
\citep[Expos\'e~IX, Theorem~6.1, p.~195]{c_sga1_2003};
\citep[Lemma~58.14.3]{BStacksFundamentalGroups}.
The splitting image is a closed subgroup of finite index
\(n=|G|\).  It defines a connected finite \'etale cover
\(P'\to P\) of degree \(n\), geometrically connected with
\(\pi_1(P'_{\overline k})=1\).  The whole section
\(\mathbf P^1_k\to P\) lifts: its fundamental-group image
is exactly this splitting subgroup.  In particular this
construction requires no finite extension of \(k\).

Its generic fibre \(B=P'_K\) is integral, being a localization
of the integral scheme \(P'\).  It is a finite \'etale cover
of \(A^\vee\) and has a \(K\)-point over zero.  Its field
of global functions is consequently \(K\), since it embeds
in \(K\) via that point.  Proper smooth base change of
\(H^0(\mathcal O)\) proves geometric connectedness.  The
pointed finite \'etale cover has a unique abelian-variety
structure making
\[
             B\longrightarrow A^\vee
\]
an isogeny of degree \(n\).  The geometric assertion is the
complete pointed-cover argument in SGA~1, Expos\'e~XI,
Section~2, pp.~221--222; uniqueness descends the group law
to \(K\).  Recall that \(p\nmid n\).

\emph{Rebuilding the model and preserving the Hodge bundle.}
One must not assume that \(H^1(P',\mathcal O_{P'})=0\).
Instead dualize the isogeny to obtain \(A\to B^\vee\).
There is a quasi-inverse whose two composites are \([n]\).
Both morphisms extend to their N\'eron models, and the induced
maps on invariant differentials compose to multiplication by
\(n\).  Because \(n\) is invertible on \(\mathbf P^1_k\),
the Hodge bundles of \(A\) and \(B^\vee\) are isomorphic.
The latter is therefore ample.  Semi-abelian reduction is
isogeny invariant
\citep[Section~1.2, Definition~1; Section~7.3, Lemma~5,
Proposition~6 and Corollary~7, pp.~180--181]{BBoschLutkebohmertRaynaud1990}.

Set \(A_{\mathrm{new}}=B^\vee\) and repeat Yuan's model
step, applying Theorem~4.5 to
\(A_{\mathrm{new}}^\vee=B\).  It constructs a model of
the identical generic abelian variety \(B\), over the same
\(K\), without another isogeny or a base extension.  Denote
this smooth projective model by \(P''\).  Its formula
\[
 R^1\psi''_*\mathcal O_{P''}
       \simeq\operatorname{Lie}(\mathcal B^\vee/\mathbf P^1_k)
\]
and the ample dual Hodge bundle give \(H^1(P'',\mathcal O)=0\).
The two smooth proper varieties \(P'\) and \(P''\) are
birational, since both model \(B/K\).  Purity and birational
invariance of the fundamental group therefore give
\[
                   \pi_1(P''_{\overline k})=1
\]
\citep[Expos\'e~X, Corollaries~3.3--3.4 and their proofs,
pp.~213--214]{c_sga1_2003}.

\emph{The surface cut.}
Apply the complete high-degree-cut construction in
Corollary~\ref{cor:B-canonical-surface-gate} to \(P''\).
When \(\dim P''=2\), first use \(P''\times\mathbf P^1\);
this preserves its geometric fundamental group, coherent
\(H^1\), and the generic Picard-zero variety \(B^\vee\).
At every subsequent smooth ample hypersurface cut, the ambient
dimension is at least three, so the map on geometric fundamental
groups is an isomorphism
\citep[Expos\'e~XII, Corollary~3.5, corrected p.~120,
original p.~152]{c_sga2_2005}.
The same Serre vanishings and finite-field Bertini conditions
preserve \(H^1(\mathcal O)=0\), injectivity on generic-fibre
\(H^1(\mathcal O)\), and smoothness of that generic fibre.
The final sufficiently high cut on a threefold makes the
canonical bundle very ample by adjunction and imposes all the
displayed numerical conditions.  Thus the resulting surface
\(Y/k\) satisfies every condition in the statement.

Let \(J\) be the Jacobian of \(Y_K\).  The tangent-space
injection makes the natural map \(B^\vee\to J\) have finite
kernel.  Hence \(B^\vee\), and therefore the isogenous
variety \(A\), is an isogeny factor of \(J\).  Tate for
\(Y\) implies BSD for \(J\), then BSD for \(A\), and then
Tate for \(X\), by the already established Yuan--Artin--Tate
chain.  This proves the equivalence over the unchanged field.
\end{proof}

The new intrinsic condition is a strict restriction of the
preceding test class in every characteristic, after a finite
constant extension if necessary.  Choose a prime \(r\geq5\)
different from \(p\), taking \(r=5\) unless \(p=5\), when
\(r=7\) suffices.  Over a finite field containing \(\mu_r\),
one may take \(k=\F_{p^{\operatorname{ord}_r(p)}}\).
The Fermat surface
\[
 F_r:\ x_0^r+x_1^r+x_2^r+x_3^r=0\quad\subset\mathbf P^3
\]
is smooth.  The cyclic group \(C_r\) acts diagonally with
weights \(0,1,2,3\).  Every nonidentity element has four
distinct eigenvalues, so its projective fixed locus is the four
coordinate points, none lying on \(F_r\).  The action is free;
its projective quotient \(Z=F_r/C_r\) is smooth and the
quotient map is finite \'etale of degree \(r\).
Lefschetz gives \(\pi_1(F_{r,\overline k})=1\), hence
\(\pi_1(Z_{\overline k})=C_r\).  The hypersurface sequence
gives \(H^1(F_r,\mathcal O)=0\), and the trace, with
\(r\) invertible, gives \(H^1(Z,\mathcal O)=0\).
Pad \(Z\) by \(\mathbf P^1\) and take a sufficiently high
smooth ample surface section.  Lefschetz preserves its
nontrivial group \(C_r\), while Serre vanishing and the
previous adjunction calculation give \(H^1(\mathcal O)=0\),
very ample canonical bundle and all the old numerical tails.
These surfaces belong to the old test class and are excluded by
the new geometric simple-connectedness condition.  This example
does not assert strictness over each individual smallest field;
the reduction theorem itself does hold over each fixed field.
The exact finite exponent checks are recorded in
\texttt{part\_b/code/fermat\_free\_action.py} and
\texttt{part\_b/output/fermat\_free\_action.txt}; the geometric
quotient and Lefschetz assertions are the proof above, not
inferences from those finite computations.
